\documentclass[11pt]{article}
\usepackage[margin=1in]{geometry}
\usepackage{amsmath,amssymb,amsthm}
\usepackage{mathtools}
\usepackage{booktabs}
\usepackage[colorlinks=true,linkcolor=blue,citecolor=blue]{hyperref}
\usepackage[capitalise,nameinlink]{cleveref}

\newtheorem{proposition}{Proposition}
\newtheorem{definition}[proposition]{Definition}
\newtheorem{remark}[proposition]{Remark}
\newtheorem{example}[proposition]{Example}

\newcommand{\Lean}{Lean}
\newcommand{\R}{\mathbb{R}}
\newcommand{\claim}[1]{\textbf{#1}}
\newcommand{\code}[1]{\texttt{#1}}

\title{A verification-first workflow for research mathematics:\\
separating machine-checked, human-proved, conditional,\\
and open claims}
\author{Fan Li\\[8pt]
{\small School of Petroleum Engineering, China University of Petroleum (East China),}\\[2pt]
{\small Qingdao 266580, China}}
\date{}

\emergencystretch=1.5em

\begin{document}
\maketitle

\begin{abstract}
A research paper in analysis makes assertions of several different strengths
within the same page: a theorem that is completely proved, a lemma whose proof
omits a standard but unpublished estimate, a conclusion that holds only under
an unproved hypothesis, a conjecture stated as a remark. When a paper is
written and revised over months, these strengths drift. We describe a workflow
that makes the strength of every assertion a machine-readable attribute,
attaches it to a specific declaration when a proof has been checked by a proof
assistant, and re-checks the whole set on every change. The workflow separates
two kinds of work: deterministic bookkeeping (which is code, with unit tests
and a self-test that deliberately breaks each gate to confirm it fires) and
mathematical judgement (which is not code, and is subjected to independent
adversarial audit).
We report a case study: an ancient axisymmetric Navier--Stokes paper with
thirty-five ledger entries, of which two were brought to \Lean~4 + Mathlib with
a clean axiom footprint, and nine were found to be blocked by missing library
theory rather than by effort. We describe the failure modes the
workflow actually caught, including two that no compiler would catch: a
formal statement that silently weakened the paper's statement at a boundary
point, and a lemma whose printed hypotheses were too weak for its printed
conclusion. We are explicit about what the workflow does not do: it does not
verify informal mathematics, it does not certify that a formal statement is the
intended one, and its audit layer is a set of prompts and adversarial reviews,
not an algorithm.
\end{abstract}

\tableofcontents

\section{Introduction}

\subsection{The problem: claim drift}

A mathematical manuscript is not a list of theorems with proofs attached. It is
a layered document in which the \emph{strength} of each assertion varies
continuously with location: the abstract says ``we prove'', Section 3 says
``we prove, assuming'', Section 5 says ``this should be standard'', and a
footnote says ``we were unable to close this step''. The reader reconstructs
the true status from context. This works, approximately, when the paper is
short and the reader is an expert in exactly this subfield.

It fails in three specific ways.

\begin{enumerate}
\item \textbf{Promotion.} A conditional statement loses its hypothesis as it
  travels from the theorem statement into the abstract, the introduction, or a
  later citation of the paper by someone else.
\item \textbf{Loss.} A step that was proved by asserting ``the standard
  maximum principle applies'' enters the main line silently, and the fact that
  the applicable notion of solution was never checked disappears.
\item \textbf{Drift.} A paper revised over months accumulates statements that
  were true of an earlier version of the argument. Nothing in the file format
  notices.
\end{enumerate}

Machine-checked mathematics addresses \emph{loss} for the steps that are
formalized, and nothing else. A proof assistant does not know what the paper
claims; it knows what the user wrote in its input language. The gap between
those two objects is where the interesting failures now occur, and it is not
made smaller by proving more theorems.

\subsection{What we describe}

We describe a workflow, implemented as a small tool (\code{mathlean}) and a
set of staged procedures, whose organizing principle is:

\begin{quote}
Every occurrence of ``we prove $X$'' in the paper must resolve to a row of a
ledger, and every row must carry one of a small number of claim states.
\end{quote}

The states are not a grading of difficulty. They are a partition by
\emph{what kind of evidence exists}:

\begin{center}
\begin{tabular}{@{}ll@{}}
\toprule
State & Evidence \\
\midrule
\claim{PROVED} & a proof assistant accepted a declaration, with a checked axiom set \\
\claim{PROVED-INFORMAL} & a complete human proof exists in the paper \\
\claim{CONDITIONAL} & a complete proof of an implication; the hypothesis is unproved \\
\claim{OPEN} & neither a proof nor a counterexample \\
\claim{REFUTED} & a counterexample, recorded \\
\claim{CITED} & attributed to another work, not proved here \\
\bottomrule
\end{tabular}
\end{center}

The ledger is a table in a Markdown file, checked into the same repository as
the manuscript. A script verifies its internal consistency, verifies that every
\claim{PROVED} row names a declaration that exists in the source, and verifies
that every named declaration has been axiom-audited.

The workflow adds two things that are, in our experience, the ones that
matter.

\textbf{First, a deterministic core.} The checking is code with unit tests, and
the checking code is itself tested by deliberately constructing broken projects
and confirming that each gate fires. A gate that has never been seen to fail is
decoration.

\textbf{Second, an adversarial audit layer.} The steps that require judgement
--- is this statement the paper's statement? does this proof use its own
conclusion? does this limit commute? --- are not handed to a single reviewer.
They are attacked by several independent reviewers with different assigned
lenses, and only ``confirmed'' findings are allowed to modify the text.

\subsection{Artifacts}

The checker, the case-study formalization, and the ledger of the two papers are
distributed as source repositories so that every claim in this document can be
re-run rather than believed:
\url{https://github.com/violetalex/set-math-ablaze}.

\subsection{What this is not}

We are explicit, because the framing invites overstatement.

\begin{itemize}
\item This is not a method for proving theorems automatically. The mathematics
  is human (or, in our runs, human-directed model output) mathematics. The
  workflow organizes evidence; it does not produce it.
\item This is not a guarantee of correctness. A \claim{PROVED} row guarantees
  that a certain formal statement follows from certain formal hypotheses under
  a checked axiom set. It does not guarantee that the formal statement is the
  intended one. That is the job of a separate audit, and the audit is
  fallible.
\item The audit layer is a protocol and a set of prompts. It has no correctness
  theorem. We report what it caught and what it missed, and the misses are part
  of the case study.
\end{itemize}

\section{The ledger}

\subsection{Format}

The ledger is a Markdown table with six columns: identifier, state, statement,
\Lean{} declaration (if any), dependencies, and notes. It is deliberately a
human-readable file rather than a serialization format, because it has to be
read and edited by the author and eventually typeset as an appendix.

Two parser details turned out to matter in practice, and both are
consequences of choosing Markdown:

\begin{itemize}
\item The identifier column is matched in full, so an identifier containing a
  lowercase suffix is silently ignored. In our first run an entry
  \code{A9a} was dropped by the parser and the ledger silently remained one
  row shorter than the author believed. The fix was to change the entry, not
  the parser; but the general lesson is that a lenient format plus a strict
  parser produces silent loss.
\item The declaration column holds one name, not a list. An entry naming two
  declarations fails the ``declaration exists'' check for the concatenation.
  This forces one row per checkable object, which is the right granularity
  anyway.
\end{itemize}

\subsection{What the ledger is checked against}

The ledger gate verifies four things:

\begin{enumerate}
\item every row's state is one of the six legal values;
\item every \claim{CONDITIONAL} row has a nonempty note identifying its
  hypothesis, and every \claim{REFUTED} row has a note containing a
  counterexample;
\item every \claim{PROVED} row names a declaration that the source actually
  contains, and that declaration appears in the axiom-audit file;
\item the dependency relation is acyclic.
\end{enumerate}

The fourth item is not bookkeeping. In our case study the most serious
mathematical error found during the audit was a circularity of exactly this
shape: a lemma $G_1$ was used to conclude that a leading-order transport term
vanishes, while the boundedness of $G_1$ was \emph{equivalent} to the vanishing
of that term. The implication was true, the proof was correct, and the
conclusion was empty. A dependency check does not detect this by itself ---
the two statements do not literally depend on each other --- but it does
force the equivalence to be written down, which is what made it visible.

\section{The pipeline}

The workflow is organized in ten stages, each with a named artifact and a pass
condition. Stages are not a waterfall: P2--P4 form a loop, and P6--P7 form
another.

\begin{center}
\begin{tabular}{@{}p{0.06\textwidth}p{0.22\textwidth}p{0.56\textwidth}@{}}
\toprule
Stage & Artifact & Pass condition \\
\midrule
P0 & \code{SPEC.md} & objects, hypotheses, conclusion with explicit
  quantifiers, and at least three falsifiability tests \\
P1 & \code{CITATIONS.md} & every reference carries an identifier and has been
  independently re-opened \\
P2 & \code{LEMMAS.md} & each lemma has statement, proof, dependencies, state \\
P3 & \code{COUNTEREXAMPLES.md} & each lemma attacked by at least one lens; no
  unexplained numerical violation \\
P4 & \code{AUDIT.md} & each finding classified \textbf{confirmed} /
  \textbf{plausible} / \textbf{refuted} \\
P5 & \code{FORMALIZATION.md} & each lemma mapped to a \Lean{} statement or
  marked out of scope with the missing theory named \\
P6 & \code{Lean/} & five gates green \\
P7 & \code{FIDELITY.md} & each formal statement compared to the specification
  and classified \textbf{match} / \textbf{weaker} / \textbf{stronger} /
  \textbf{different} \\
P8 & \code{CLAIMS.md} & every assertion in the paper resolves to a row \\
P9 & \code{FINAL.md} & the terminal checklist \\
\bottomrule
\end{tabular}
\end{center}

Two design decisions in this table are worth isolating.

\textbf{P0 forbids generalization.} The most common way for a statement-locking
stage to fail is to ``strengthen'' the problem. In a rehearsal run on the
elementary inequality $2ab \le a^2+b^2$, the specification stage generalized
the statement to arbitrary linearly ordered commutative rings. Every
downstream attack then explored zero divisors, and no agent touched the actual
inequality. The rule we adopted is that generalizations may appear only in an
``out of scope'' section, never in the hypotheses or the conclusion.

\textbf{P5 is an estimate, not a judgement of truth.} The formalization
mapping classifies each lemma as \emph{easy}, \emph{medium}, \emph{hard}, or
\emph{not feasible}, where the classes are statements about the state of the
library, not about the mathematics. A false lemma can be \emph{medium} and a
true one \emph{not feasible}. Conflating the two is a category error, and we
have seen it made.

\section{The deterministic core}

\subsection{What is code}

Everything that can be decided mechanically is implemented in a dependency-free
Python package of roughly two thousand lines, with 86 unit tests and 12
self-test scenarios:

\begin{itemize}
\item scaffolding a project and vendoring the checker into it;
\item comment-aware scanning for \code{sorry} and \code{admit};
\item running the axiom audit and parsing its output;
\item parsing and checking the ledger;
\item checking that the stage artifacts exist;
\item reporting, with an actionable fix, why the environment is broken.
\end{itemize}

The comment-aware scanner deserves a remark, because our first implementation
was a regular expression. It matched the word \code{sorry} inside a
documentation string that said ``\code{sorry} is not allowed'', and the gate
went red on a correct project. Real \Lean{} source has nestable block comments,
line comments, and string literals containing dashes. The scanner now walks the
source character by character and preserves line numbering, because a scanner
whose reported line numbers drift is a scanner nobody trusts.

\subsection{The five gates}

\begin{center}
\begin{tabular}{@{}ll@{}}
\toprule
Gate & Question \\
\midrule
\code{build} & does the library compile? \\
\code{sorry} & is there an admitted hole? \\
\code{axioms} & does each public declaration depend only on permitted axioms? \\
\code{ledger} & does every claim resolve, and is the dependency graph acyclic? \\
\code{docs} & are the stage artifacts present? \\
\bottomrule
\end{tabular}
\end{center}

The first three gates concern whether something was proved. The fourth
concerns whether the thing proved is the thing claimed. The fifth is
hygiene.

The axiom gate is stronger than ``the build succeeded'' in a way that
matters. A \Lean{} file can compile while depending on \code{sorryAx} (via an
admitted lemma) or on a user-declared \code{axiom}. Neither appears in the
source of the theorem being checked. The gate therefore runs
\code{\#print axioms} on every public declaration and compares the transitive
closure against a whitelist, by default \{\code{propext},
\code{Classical.choice}, \code{Quot.sound}\}. Custom axioms are permitted
only if they are registered in the ledger with a description of what they
correspond to in the paper.

\subsection{Testing the tests}

A gate is only evidence if it has been observed to fail on the thing it
claims to catch. The self-test constructs a temporary project and breaks it in
twelve ways, one per scenario: a genuine \code{sorry}; the word
\code{sorry} inside a comment; a \claim{PROVED} row naming a
nonexistent declaration; a \claim{PROVED} row with no axiom audit; a
\claim{CONDITIONAL} row with no hypothesis; a \claim{REFUTED} row with no
counterexample; an illegal state; a dependency cycle; an unregistered custom
axiom; and a registered custom axiom that must be allowed. Each scenario
asserts both the verdict and the reason.

The self-test found real defects in the checker, including a template project
that passed the documentation gate because the scaffolding generated the
artifact files that the gate was looking for. The gate was checking for
existence, and the scaffolding guaranteed existence. It now checks that the
artifact has been edited since generation.

\section{The judgement layer}

\subsection{Why independent review}

The errors that a formalization does not catch are not arithmetic errors. They
are errors of the following kinds, all drawn from our case study:

\begin{itemize}
\item a formal statement that is weaker than the paper's at a boundary point;
\item a formal statement that is stronger than the paper's, because a
  hypothesis was dropped;
\item a proof that is correct but whose conclusion was already assumed;
\item a printed lemma whose hypotheses do not imply its printed conclusion;
\item a quantitative claim whose stated range of validity is wrong.
\end{itemize}

Each of these is invisible to a compiler and visible to a reader who is
looking for it. The workflow therefore assigns lenses. Attackers are given the
specification, the paper, and the formal source, and are told to falsify; they
are not told the author's confidence.

\subsection{Adjudication}

Findings carry one of three verdicts. \textbf{confirmed} means the reviewer
can exhibit the specific error. \textbf{plausible} means there is a concrete
doubt that has not been closed. \textbf{refuted} means the reviewer's own
objection was wrong. Only \textbf{confirmed} may modify the text;
\textbf{plausible} goes to a backlog. The loop terminates when two consecutive
rounds produce no new \textbf{confirmed} finding.

Reviewers disagree, and the disagreements are recorded rather than resolved by
majority. In our run, two auditors made opposite claims about whether a
simplified quantity made a hypothesis stronger or weaker; the record preserves
both statements, and the resolution cites the direction of an inequality. The
value of keeping the disagreement is that it shows which claims are
load-bearing and which were asserted casually.

\subsection{The cost of judgement}

Adversarial review is the expensive part, and the budget must be explicit. Our
configuration caps the number of reviewer agents per round and the number of
findings adjudicated, sorts findings by severity before spending reviewers on
them, and reports what was skipped. The failure mode to avoid is silent
truncation: a run that reviewed the first six findings and stopped reads
exactly like a run that found six findings.

\section{Case study: an ancient axisymmetric Navier--Stokes paper}

\subsection{The project}

The subject is a paper on ancient (global backward in time) axisymmetric
solutions of the three-dimensional incompressible Navier--Stokes equations with
swirl, in the scaling-critical Koch--\allowbreak Nadirashvili--\allowbreak Seregin--\allowbreak \v{S}ver\'ak class
at the endpoint $\gamma=0$ \cite{KNSS}. The paper proves a conditional rigidity
statement and reduces the remaining case to a sufficient condition on
$\R^5$; it explicitly does not claim an unconditional theorem at the endpoint.
A companion paper treats a two-dimensional lemma that closes the last step.

The manuscript was audited and formalized as part of its preparation, not
after publication. That timing is the point: the workflow found errors while
they were still cheap.

\subsection{What was mapped}

Thirty-one formal statements were extracted from the two papers and classified
by formalization cost. After deduplicating three restatements, twenty-eight
independent results remain.

\begin{center}
\begin{tabular}{@{}lrr@{}}
\toprule
Cost & before dedup & after dedup \\
\midrule
easy & 0 & 0 \\
medium & 8 & 8 \\
hard & 14 & 12 \\
not feasible & 9 & 8 \\
\bottomrule
\end{tabular}
\end{center}

Two observations from the mapping.

\textbf{There was nothing easy.} The cheapest results in the batch --- an
integral estimate, a weight trade-off computation, and a pointwise identity ---
each required between two and six hundred lines of \Lean. The reason is not
that the mathematics is deep; it is that the ambient vocabulary (cylindrical
coordinates, a weighted Laplacian on $\R^5$, axis regularity) does not exist in
the library and must be built before the first line of the intended argument
can be written.

\textbf{The obstruction is a single missing theory.} Nine results were
classified out of scope. Five of them, including the main theorem, the
weighted maximum principle and the dual-density criterion, fail at the same
place: the library has no Feynman--Kac representation, no parabolic maximum
principle, and no martingale problem for bounded measurable drift. Two more
fail at the adjacent probabilistic layer (a stochastic-analysis lemma, and a
lemma whose core is a one-dimensional Bessel hitting probability), one on the
$n$-dimensional theory of harmonic functions, and one is an external theorem
that can only be entered as a citation. This is the most useful output of the mapping, and it
is a statement about \emph{the library}, not about the mathematics: the same
theorem might be formalizable tomorrow, and the estimate tells you where to
look.

\subsection{The formalized result}

One result was brought through to a machine-checked statement: the final
integration step of the uniform bound on the lifted swirl $\Phi = \Gamma/\rho^2$,
where $\Gamma = \rho u^\theta$ and $R = 1+|x|+\sqrt{-t}$.

The paper's proof ends with
\begin{equation}
|\Gamma(\rho,\zeta,t)| \le \int_0^\rho \frac{C_1 r \, dr}{R(r,\zeta,t)}
\le \frac{C_1 \rho^2}{2(1+\sqrt{-t})},
\end{equation}
obtained from $\Gamma(0,\zeta,t)=0$ by the fundamental theorem of calculus.
The formalization covers this step and its corollary, in which the bound is
divided by $\rho^2$ and the resulting function of $t$ is shown to tend to zero
as $t \to -\infty$.

The step from the paper's hypotheses (H1)--(H3) to the pointwise bound
$|\partial_\rho \Gamma| \le C_1\rho/R$ is the identity
$\partial_\rho \Gamma = \rho\,\omega^\zeta$ together with $|\omega^\zeta| \le
|\omega| \le C_1/R$, i.e.\ a curl computation in cylindrical coordinates for a
vector field on $\R^3$. This is not formalized; it enters the theorem as a
hypothesis. The ledger row for the full statement therefore remains
\claim{PROVED-INFORMAL}, and the two \claim{PROVED} rows cover its last step
and its corollary. The distinction is visible in the repository, in the ledger,
and in the paper's own statement of the formalization's scope.

\subsection{Three failures the workflow caught}

\paragraph{A boundary point that silently weakened the statement.}
The lifted variable is $\Phi = \Gamma/\rho^2$. In \Lean, division by zero is
total and $x/0 = 0$. Writing the theorem with $\Phi$ defined by the literal term
$\Phi(\rho,\zeta,t) := \Gamma/\rho^2$ therefore makes the statement at
$\rho = 0$ read $0 \le C_1/(2(1+\sqrt{-t}))$, which is true and says nothing. The paper's $\Phi$ is a smooth lift $g(|y|^2,y_5,t)$ with
$\Phi(0,\zeta,t) = g(0,\zeta,t)$ in general nonzero. The formal statement had
been \emph{proved} and was weaker than the paper's at exactly one point.

The independent audit found this; the compiler did not, and could not. The fix
was to define the extension explicitly, with the limiting value as a parameter,
and to carry the existence of the limit as a hypothesis. The episode is the
clearest illustration we have of why P6 and P7 are separate stages: a green
build is not evidence that the right theorem was stated.

\paragraph{A hypothesis that did not imply its conclusion.}
A corollary in the paper derived a bound $h \le C\rho^2$ on $(0,\infty)$ from
the differential inequality $h'' \le h'/\rho$, together with a condition on a
power weight. The inequality is true on $[\rho_0,\infty)$ for any
$\rho_0>0$, but false on $(0,\infty)$: the function $h = 1-\rho^2$ satisfies
the inequality and is not $O(\rho^2)$ at the origin. The corollary needed a
locality hypothesis that was not printed. This was found by a numerical probe
and confirmed in \Lean, both before the paper is submitted.

\paragraph{A dependency that was an equivalence.}
A two-scale reduction argued that a certain bounded correction $G_1$ forces a
leading-order transport coefficient to vanish. The argument is valid. What the
audit noticed is that the boundedness of $G_1$ is equivalent, in the class
considered, to the vanishing of that coefficient. The lemma is therefore
circular in the sense that it assumes what it concludes, and the paper now
records the reduction as formal rather than as progress. The published version
does not claim otherwise.

\subsection{Cost}

For the formalized lemma, the accounting was:

\begin{center}
\begin{tabular}{@{}lrr@{}}
\toprule
Stage & agents & tokens \\
\midrule
cost ranking over 31 statements & 7 & 597{,}791 \\
independent audit of one lemma & 4 & 241{,}500 \\
\bottomrule
\end{tabular}
\end{center}

We report these because the alternative --- describing the workflow as cheap
--- would be false. The audit was worth its cost: it found the boundary-point
weakening, which is a defect in the formalization that would otherwise have
been published as a success.

\section{Observations from operation}

The workflow is a program that runs other programs and, in our setup, model
instances. Several of its failures were in the orchestration rather than in
the mathematics, and we record them because they are reproducible and because
they shaped the design.

\paragraph{Structured verdicts.} An early version judged success by whether a
stage's returned text began with ``OK''. A stage that correctly reported
``declaration exists, compilation succeeded, verification follows'' was scored
as a failure, blocking eleven downstream lemmas. Verdicts are now structured
values, and text is for humans.

\paragraph{Todo lists keyed by identity, not by value.} Removing completed
items by object identity failed once the objects crossed a dispatch boundary:
118 dispatches went to four lemmas that had already succeeded, while the
remaining ten were never attempted. Keying the todo list by identifier removed
the failure.

\paragraph{Naming the deadlock.} When no lemma is ready, the correct behavior
is to report which lemma is blocked by which, and stop. Retrying a random
subset produces activity and no progress; in one rehearsal, the majority of
139 dispatches fell into this category.

\paragraph{Silent skip is a lie.} A stage that returns nothing must be
distinguishable from a stage that found nothing. Every bounded stage reports
its bound.

\section{Related work}

Proof assistants have been used for large-scale formalization projects for
over a decade \cite{Hales,OddOrder}, and the mathematical library
\cite{mathlib} behind \Lean~4 \cite{Lean4} now covers substantial parts of
undergraduate and graduate mathematics. Those projects establish that
formalization at scale is possible.

Our concern is different and complementary. A large formalization project has
a specification: the statement to be formalized is fixed by a committee or by a
prior informal text, and the engineering problem is to reach it. A research
paper has no such specification; the statement is being discovered while it is
being proved, and the artifact that must be maintained is not the proof but the
\emph{map} from claims to evidence. The ledger is that map.

The workflow differs from a purely formal approach in one further respect: it
treats the informal paper as the primary object and the formalization as one
kind of evidence among several. A paper whose main theorem is not formalizable
in the current state of the library is not a failed formalization project; it
is a paper with an accurate statement about what is and is not checked.

\section{Limitations}

\begin{enumerate}
\item \textbf{The audit is not an algorithm.} Independent review with assigned
  lenses found real errors in our case study, and the same protocol missed
  others that we do not know about. There is no completeness claim.
\item \textbf{Fidelity is judged, not certified.} The comparison between a
  formal statement and an informal one is a human judgement recorded in a
  document. A mis-specified formalization can pass every gate.
\item \textbf{The cost estimates are local.} ``Not feasible'' means not
  feasible with the library as it stands, for an author with the author's time.
  It is not a statement about difficulty.
\item \textbf{The case study is one project.} The failure modes we report are
  those that occurred in a paper about a conditional Liouville theorem. Other
  fields will produce others.
\item \textbf{Model involvement.} Judgement stages in our runs were executed
  with model assistance, directed and reviewed by the author. The protocol is
  designed to make that safe --- every output lands in a file, every file is
  gated, every gate is self-tested --- but the protocol does not remove the
  dependence on the author's reading.
\end{enumerate}

\section{Conclusion}

The contribution we claim is narrow and, we think, transferable: make the
strength of every assertion a first-class, checked attribute, and make the
checks themselves testable. The result is a paper that states, in a form a
script can verify, which of its claims are machine-checked, which are proved on
paper, which depend on an unproved hypothesis, and which are open.

For the case study, that meant two declarations with a clean axiom footprint, a
public record of the thirty-three claims that are not machine-checked, and a
list of nine results blocked by one missing theory rather than by one missing
idea.

\appendix

\section{Gate definitions}

\begin{description}
\item[\textbf{build}] \code{lake build} exits zero.
\item[\textbf{sorry}] the comment-stripped source contains no \code{sorry}
  or \code{admit}.
\item[\textbf{axioms}] \code{Audit.lean} compiles, and for every declaration
  named in it, the transitive axiom set is contained in the configured
  whitelist. A custom axiom not registered in the ledger is an error.
\item[\textbf{ledger}] the ledger parses; states are legal; \claim{CONDITIONAL}
  rows carry hypotheses; \claim{REFUTED} rows carry counterexamples;
  \claim{PROVED} rows name existing declarations that are axiom-audited; the
  dependency graph is acyclic.
\item[\textbf{docs}] the stage artifacts exist and have been modified since
  scaffolding.
\end{description}

\section{An excerpt of the ledger}

\begin{center}
\begin{tabular}{@{}p{0.05\textwidth}p{0.20\textwidth}p{0.27\textwidth}p{0.26\textwidth}@{}}
\toprule
ID & State & Claim & Declaration \\
\midrule
A2 & \claim{CONDITIONAL} & $\Gamma \equiv 0$ under $b^\rho \ge 0$ & --- \\
A4 & \claim{CONDITIONAL} & $\Gamma \equiv 0$ if the scale-free condition
  $(\star)$ admits a bounded supersolution & --- \\
A9 & \claim{PROVED-INFORMAL} & uniform backward bound on $\Phi$ & --- \\
A10 & \claim{PROVED-INFORMAL} & among power weights, only $\rho^2$ is
  compatible & --- \\
A13 & \claim{PROVED-INFORMAL} & averaged necessary condition & --- \\
A27 & \claim{OPEN} & the $\gamma=0$ endpoint in general & --- \\
A34 & \claim{PROVED} & integration step of A9, for all $\rho \ge 0$ &
  \code{PhiExt\_bound} \\
A35 & \claim{PROVED} & corollary: $\Phi \to 0$ uniformly in the past &
  \code{PhiExt\_bound\_tendsto} \\
\bottomrule
\end{tabular}
\end{center}

Of the thirty-five rows, two are \claim{PROVED}, nineteen
\claim{PROVED-INFORMAL}, nine \claim{CONDITIONAL}, two \claim{OPEN}, and three
\claim{CITED}.

\begin{thebibliography}{99}

\bibitem{KNSS}
G.~Koch, N.~Nadirashvili, G.~Seregin, and V.~\v{S}ver\'ak.
\newblock Liouville theorems for the Navier--Stokes equations and applications.
\newblock \emph{Acta Mathematica} \textbf{203} (2009), 83--105.

\bibitem{Lean4}
L.~de Moura and S.~Ullrich.
\newblock The Lean 4 theorem prover and programming language.
\newblock \emph{Automated Deduction -- CADE 28}, Lecture Notes in Computer
  Science 12699, Springer, 2021, 625--635.

\bibitem{mathlib}
The mathlib Community.
\newblock The Lean mathematical library.
\newblock \emph{Proceedings of the 9th ACM SIGPLAN International Conference on
  Certified Programs and Proofs (CPP 2020)}, 2020, 367--381.

\bibitem{Hales}
T.~Hales et al.
\newblock A formal proof of the Kepler conjecture.
\newblock \emph{Forum of Mathematics, Pi} \textbf{5} (2017), e2.

\bibitem{OddOrder}
G.~Gonthier et al.
\newblock A machine-checked proof of the odd order theorem.
\newblock \emph{Interactive Theorem Proving (ITP 2013)}, Lecture Notes in
  Computer Science 7998, Springer, 2013, 163--179.

\end{thebibliography}

\end{document}
